\documentclass[a4paper,12pt]{article}

% ===========================================================================
%  Лабораторная работа 3.5.1 «Изучение плазмы газового разряда в неоне»
%
%  ЗАГОТОВКА ОТЧЁТА (до измерений). Теория, методика и структура разделов
%  готовы; места под данные помечены \qq (число) и \fillme (текст).
%
%  Сборка:  cd 3.5.1/report && pdflatex main.tex  (или latexmk -pdf main.tex)
%  Таблицы и графики читаются из ../data/processed/*.csv, которые готовит
%  ../scripts/process_data.py. Пока файлов нет, вместо таблиц выводятся
%  пустые бланки (их удобно распечатать и заполнять во время работы),
%  вместо графиков -- рамки-заглушки.
% ===========================================================================

% ============================================================
% Общий преамбул для отчётов по лабораторным работам (labs_lpr)
% Подключается из report/main.tex командой:
%   % ============================================================
% Общий преамбул для отчётов по лабораторным работам (labs_lpr)
% Подключается из report/main.tex командой:
%   % ============================================================
% Общий преамбул для отчётов по лабораторным работам (labs_lpr)
% Подключается из report/main.tex командой:
%   \input{../../common/latex_template/preamble.tex}
% ============================================================

% Поддержка русского языка и кодировки
\usepackage[T2A]{fontenc}
\usepackage{lmodern}
\usepackage[utf8]{inputenc}
\usepackage[russian]{babel}

% Математические пакеты
\usepackage{amsmath, amssymb}

% Настройка полей
\usepackage{geometry}
\geometry{top=2cm, bottom=2cm, left=2.5cm, right=1.5cm}

% Колонтитулы
\usepackage{fancyhdr}
\usepackage{graphicx}

% Современные таблицы
\usepackage{tabularray}

% Работа с csv файлами
\usepackage{csvsimple-l3}
\usepackage{tikz}

% Графики
\usepackage{pgfplots}
\pgfplotsset{compat=1.18}

% Ссылки
% \usepackage{hyperref} % гипперссылки, нужно оформление вместо красного квадрата
\usepackage[russian]{cleveref} % умные ссылки, нужно переоформление для формул, рисунков и тд (вместо ф-л что-то нужно придумать)

% ============================================================

% Поддержка русского языка и кодировки
\usepackage[T2A]{fontenc}
\usepackage{lmodern}
\usepackage[utf8]{inputenc}
\usepackage[russian]{babel}

% Математические пакеты
\usepackage{amsmath, amssymb}

% Настройка полей
\usepackage{geometry}
\geometry{top=2cm, bottom=2cm, left=2.5cm, right=1.5cm}

% Колонтитулы
\usepackage{fancyhdr}
\usepackage{graphicx}

% Современные таблицы
\usepackage{tabularray}

% Работа с csv файлами
\usepackage{csvsimple-l3}
\usepackage{tikz}

% Графики
\usepackage{pgfplots}
\pgfplotsset{compat=1.18}

% Ссылки
% \usepackage{hyperref} % гипперссылки, нужно оформление вместо красного квадрата
\usepackage[russian]{cleveref} % умные ссылки, нужно переоформление для формул, рисунков и тд (вместо ф-л что-то нужно придумать)

% ============================================================

% Поддержка русского языка и кодировки
\usepackage[T2A]{fontenc}
\usepackage{lmodern}
\usepackage[utf8]{inputenc}
\usepackage[russian]{babel}

% Математические пакеты
\usepackage{amsmath, amssymb}

% Настройка полей
\usepackage{geometry}
\geometry{top=2cm, bottom=2cm, left=2.5cm, right=1.5cm}

% Колонтитулы
\usepackage{fancyhdr}
\usepackage{graphicx}

% Современные таблицы
\usepackage{tabularray}

% Работа с csv файлами
\usepackage{csvsimple-l3}
\usepackage{tikz}

% Графики
\usepackage{pgfplots}
\pgfplotsset{compat=1.18}

% Ссылки
% \usepackage{hyperref} % гипперссылки, нужно оформление вместо красного квадрата
\usepackage[russian]{cleveref} % умные ссылки, нужно переоформление для формул, рисунков и тд (вместо ф-л что-то нужно придумать)


% Пакеты, нужные именно этой работе
\usepackage{subcaption}
\usepackage{circuitikz}

% --- Служебные макросы ------------------------------------------------------
% Короткий пропуск (число/слово, которое надо вписать после измерений)
\newcommand{\qq}{\textbf{???}}
% Пропуск, требующий осмысленного текста
\newcommand{\fillme}{\textbf{?? НЕОБХОДИМО ЗАПОЛНИТЬ ??}}
% Имя файла/пути моноширинным шрифтом
\newcommand{\filename}[1]{\texttt{\detokenize{#1}}}
% Рисунок из report/img/; если файла ещё нет -- рамка-заглушка
\newcommand{\labimage}[2]{%
    \IfFileExists{img/#1}{%
        \includegraphics[width=#2\linewidth]{img/#1}%
    }{%
        \fbox{\parbox[c][0.6\dimexpr#2\linewidth\relax][c]{%
                \dimexpr#2\linewidth-2\fboxsep-2\fboxrule\relax}{%
            \centering\small\filename{img/#1}\\[2mm]\fillme}}%
    }%
}
% Путь к обработанным данным
\newcommand{\proc}[1]{../data/processed/#1}

% Числа из CSV: десятичная запятая; пустая ячейка (NaN) -> «--»
\pgfkeys{/pgf/number format/.cd, use comma, 1000 sep={\,}}
\pgfkeys{/pgf/number format/csvnum/.style={precision=6, std=-3:5}}
\ExplSyntaxOn
\NewDocumentCommand{\pnum}{m}{
    \tl_if_blank:eTF {#1} {--} { \pgfmathprintnumber[csvnum]{#1} }
}
% Пустой бланк таблицы: #1 -- colspec, #2 -- строка заголовка,
% #3 -- число столбцов, #4 -- число пустых строк
\NewDocumentCommand{\BlankTable}{mmmm}{
    \tl_set:Ne \l_tmpa_tl {
        \exp_not:N \begin{tblr}
            \exp_not:n {{colspec={#1}, hlines, vlines,
                         row{1}={font=\bfseries}, rows={ht=6.5mm}}}
        \exp_not:n {#2} \exp_not:N \\
        \prg_replicate:nn {#4} { \prg_replicate:nn {#3 - 1} { & } \exp_not:N \\ }
        \exp_not:N \end{tblr}
    }
    \tl_use:N \l_tmpa_tl
}
\ExplSyntaxOff

% Колонтитулы (как в остальных работах репозитория)
\pagestyle{fancy}
\fancyhf{}
\fancyhead[L]{%
    \IfFileExists{../../common/latex_template/miptlogo.png}{%
        \includegraphics[height=1cm]{../../common/latex_template/miptlogo.png}%
    }{%
        \textbf{[ЛОГОТИП МФТИ]}%
    }%
}
\fancyhead[R]{\textbf{Емельянов ДА, Байдин ФМ}}
\fancyfoot[C]{\thepage}
\renewcommand{\headrulewidth}{0.4pt}
\setlength{\headheight}{33pt}

\begin{document}

% --- ТИТУЛЬНЫЙ ЛИСТ ---
\begin{titlepage}
\begin{center}
    {\large МОСКОВСКИЙ ФИЗИКО-ТЕХНИЧЕСКИЙ ИНСТИТУТ}
\end{center}
\begin{center}
    {\large Физтех-школа физики и исследований им. Ландау\break ЛФИ}
\end{center}

    \vspace{3cm}

\vspace{0.15cm}
{
\begin{center}
    \huge \bf Лабораторная работа №3.5.1 по курсу Общая физика\\
    \vspace{0.4cm}
    \LARGE Изучение плазмы газового разряда в неоне.
\end{center}
}
\vspace{1.5cm}
\begin{flushright}
Емельянов Дмитрий Алексеевич\\
Байдин Фаддей Михайлович\\
\vspace{0.4cm}
Б02-502
\end{flushright}
\vspace{10cm}
\begin{center}
    \large \qq{} 2026, г. Долгопрудный
\end{center}
\end{titlepage}
% ----------------------

\section{Аннотация}
\textbf{Цель работы:} изучение вольт-амперной характеристики тлеющего
разряда в неоне; изучение свойств плазмы положительного столба разряда
методом двойного зонда.

В работе измеряется вольт-амперная характеристика (ВАХ) тлеющего разряда
$I_{\text{р}}(U_{\text{р}})$ и определяется его дифференциальное
сопротивление. При нескольких значениях тока разряда снимаются
вольт-амперные характеристики двойного зонда, по которым находятся ионный
ток насыщения и наклон характеристики в нуле. По ним рассчитываются
температура $T_e$ и концентрация $n_e$ электронов, плазменная частота
$\omega_p$, электронная поляризационная длина $r_{De}$, дебаевский радиус
экранирования $r_D$, число частиц в дебаевской сфере $N_D$ и степень
ионизации плазмы $\alpha$, а также исследуются зависимости $T_e(I_{\text{р}})$
и $n_e(I_{\text{р}})$.

\section{Теоретические сведения}
В теоретическом разделе, как и в методическом пособии, используется система
СГС; расчётные формулы в СИ приведены в \cref{sec:si}.

\subsection{Плазма. Плазменная частота}
\emph{Плазмой} называют ионизованный газ, поведение заряженных частиц в
котором носит коллективный характер. Плазма \emph{квазинейтральна}: средние
концентрации электронов и однозарядных ионов равны, $n_e = n_i = n$.

Если сместить все электроны некоторого объёма плазмы на расстояние
$\Delta x$ относительно (неподвижных, тяжёлых) ионов, возникнет поле
$E = 4\pi n e\,\Delta x$, возвращающее электроны назад. Уравнение движения
электрона $\ddot{\Delta x} = -\frac{4\pi n e^2}{m_e}\Delta x$ описывает
гармонические (ленгмюровские) колебания с \emph{плазменной частотой}
\begin{equation}
    \label{eq:wp}
    \omega_p = \sqrt{\frac{4\pi n_e e^2}{m_e}}.
\end{equation}
Она задаёт характерный временной масштаб отклика плазмы на нарушение
квазинейтральности.

\subsection{Дебаевский радиус. Идеальность плазмы}
Амплитуда плазменных колебаний, возбуждаемых тепловым движением электронов,
$r \sim \bar v_e/\omega_p$, определяет \emph{дебаевский радиус}
\begin{equation}
    \label{eq:rD}
    r_D = \sqrt{\frac{k_{\text{Б}} T}{4\pi n e^2}}.
\end{equation}
На расстояниях порядка $r_D$ поле пробного заряда в плазме экранируется:
$\varphi(r) = \frac{q}{r}e^{-r/r_D}$. В неравновесной плазме
($T_e \neq T_i$) вводят две длины
\begin{equation}
    \label{eq:rDe_rDi}
    r_{De} = \sqrt{\frac{k_{\text{Б}} T_e}{4\pi n_e e^2}},
    \qquad
    r_{Di} = \sqrt{\frac{k_{\text{Б}} T_i}{4\pi n_i e^2}},
\end{equation}
а радиус экранирования равен
\begin{equation}
    \label{eq:rD_full}
    r_D = \left(r_{De}^{-2} + r_{Di}^{-2}\right)^{-1/2}
    \approx r_{Di} \quad\text{при}\quad T_e \gg T_i.
\end{equation}
Электронную длину $r_{De}$ называют также \emph{поляризационной}: она
определяет масштаб, на котором квазинейтральность нарушается тепловыми
флуктуациями электронов. Ионизованный газ является плазмой, если
$r_D \ll a$, где $a$ --- характерный размер системы.

Отношение тепловой энергии частиц к энергии их кулоновского взаимодействия
по порядку величины равно числу частиц в дебаевской сфере
\begin{equation}
    \label{eq:ND}
    N_D = \frac43\pi n_i r_D^3.
\end{equation}
Плазма \emph{идеальна}, если $N_D \gg 1$.

В тлеющем разряде плазма существенно неравновесна: электроны,
ускоряемые полем, имеют $T_e \sim 10^4$~К, а ионы, эффективно обменивающиеся
энергией с атомами газа, --- $T_i \sim 300$~К.

\subsection{Зондовый метод. Двойной зонд}
Проводник (зонд), внесённый в плазму, заряжается до отрицательного
\emph{плавающего потенциала} $U_f$, при котором электронный и ионный токи на
него уравновешиваются ($\bar v_e \gg \bar v_i$):
\begin{equation}
    \label{eq:Uf}
    U_f = -\frac12\frac{k_{\text{Б}} T_e}{e}\ln\frac{T_e m_i}{T_i m_e}.
\end{equation}
Ионный ток насыщения на зонд площади $S$ определяется не тепловой
скоростью ионов, а их ускорением в приповерхностном слое; для расчётов
используется полуэмпирическая формула Бома
\begin{equation}
    \label{eq:bohm}
    I_{i\text{н}} \approx 0{,}4\, n_i e S\sqrt{\frac{2k_{\text{Б}} T_e}{m_i}}.
\end{equation}

\emph{Двойной зонд} --- два одинаковых близко расположенных зонда, между
которыми создаётся разность потенциалов $U$, $|U| \ll |U_f|$. Оба зонда
находятся вблизи плавающего потенциала, на ионной ветви ВАХ, и ионные токи
на них равны $I_{i\text{н}}$. Ток через систему определяется разностью
электронных токов; с учётом больцмановского распределения электронов и
условия $I_1 = -I_2 = I$ получается
\begin{equation}
    \label{eq:tanh}
    I = I_{i\text{н}}\,\operatorname{th}\frac{eU}{2k_{\text{Б}} T_e}.
\end{equation}
При $|U| \to \infty$ ток стремится к $\pm I_{i\text{н}}$; реальная
характеристика имеет наклонные асимптоты из-за ускорения частиц
приложенным полем, поэтому $I_{i\text{н}}$ находится по пересечению
асимптот с осью $U = 0$. Дифференцируя \eqref{eq:tanh} в нуле, получаем
формулу для электронной температуры
\begin{equation}
    \label{eq:Te}
    k_{\text{Б}} T_e = \frac12\,\frac{e\, I_{i\text{н}}}
                           {\left.\dfrac{dI}{dU}\right|_{U=0}}.
\end{equation}
По известным $T_e$ и $I_{i\text{н}}$ из \eqref{eq:bohm} находится
концентрация $n_e = n_i$. Площадь зонда в пренебрежении дебаевским слоем
$S \approx \pi d l$.

\subsection{Вольт-амперная характеристика тлеющего разряда}
Самостоятельный разряд в газе зажигается при напряжении зажигания
$U_{\text{заж}}$, когда выполнен критерий Таунсенда: каждый электрон,
вышедший из катода, за счёт ударной ионизации в газе и вторичной эмиссии
ионами с катода воспроизводит хотя бы один новый электрон. Гаснет разряд
при меньшем напряжении, чем зажигается (гистерезис).

ВАХ газового разряда в широком диапазоне токов
(рис.~6 Приложения к разделу V методического пособия) содержит участки:
тёмный таунсендовский разряд; поднормальный тлеющий разряд (падающий
участок); \emph{нормальный} тлеющий разряд, при котором напряжение почти не
зависит от тока (растёт лишь площадь катодного пятна при постоянной
плотности тока); \emph{аномальный} тлеющий разряд (напряжение растёт с
током); переход к дуговому разряду. На падающих участках
дифференциальное сопротивление
\begin{equation}
    \label{eq:Rdiff}
    R_{\text{диф}} = \frac{dU}{dI}
\end{equation}
отрицательно. Для устойчивого горения разряда последовательно с ним
включают балластный резистор $R_{\text{б}}$: рабочая точка лежит на
пересечении ВАХ с нагрузочной прямой $I = (\mathcal{E} - U)/R_{\text{б}}$,
а условие устойчивости имеет вид $R_{\text{диф}} + R_{\text{б}} > 0$.

\subsection{Расчётные формулы в СИ}
\label{sec:si}
Обработка (скрипт \filename{scripts/process_data.py}) ведётся в СИ:
\begin{gather}
    k_{\text{Б}} T_e\,[\text{эВ}] = \frac{I_{i\text{н}}}{2\,dI/dU},
    \qquad
    n_e = \frac{I_{i\text{н}}}{0{,}4\, e S}\sqrt{\frac{m_i}{2 k_{\text{Б}} T_e}},
    \qquad S = \pi d l,
    \label{eq:si1}\\
    \omega_p = \sqrt{\frac{n_e e^2}{\varepsilon_0 m_e}},
    \qquad
    r_{De} = \sqrt{\frac{\varepsilon_0 k_{\text{Б}} T_e}{n_e e^2}},
    \qquad
    r_{Di} = \sqrt{\frac{\varepsilon_0 k_{\text{Б}} T_i}{n_e e^2}},
    \label{eq:si2}\\
    N_D = \frac43\pi n_e r_D^3,
    \qquad
    \alpha = \frac{n_e}{n_0},\quad n_0 = \frac{P}{k_{\text{Б}} T},
    \label{eq:si3}
\end{gather}
где $m_i = 22\,\text{а.е.м.}$ (изотоп $^{22}$Ne), $T_i \approx T = 300$~К,
$P \approx 2$~торр.

\section{Используемое оборудование}
\begin{itemize}
    \item стеклянная газоразрядная трубка с холодным полым катодом, тремя
          анодами и геттерным узлом, наполненная неоном $^{22}$Ne при
          давлении $2$~мм рт.\,ст.;
    \item двойной зонд из молибденовой проволоки
          (диаметр $d = \qq$~мм, длина $l = \qq$~мм);
    \item высоковольтный источник питания (ВИП) до нескольких киловольт;
    \item балластный резистор $R_{\text{б}} \approx 500$~кОм
          ($R_{\text{б}} = \qq$~кОм);
    \item высокоомный делитель напряжения $(R_1 + R_2)/R_2 = \qq$;
    \item источник постоянного тока с дискретным переключателем
          выходного напряжения и потенциометр $R$;
    \item миллиамперметр $A_1$ (ток разряда), вольтметр $V_1$
          (напряжение на трубке), микроамперметр $A_2$ (ток зонда),
          вольтметр $V_2$ (напряжение на зондах);
    \item переключатели П1 (выбор анода) и П2 (полярность напряжения на
          зондах).
\end{itemize}

\section{Методика измерений}
Схема установки приведена на \cref{fig:scheme}. Катод и один из анодов
(I или II) через переключатель П1 и балластный резистор $R_{\text{б}}$
подключаются к ВИП. Ток разряда измеряется миллиамперметром $A_1$,
напряжение на трубке --- вольтметром $V_1$ через делитель с коэффициентом
$(R_1+R_2)/R_2$.

\begin{figure}[!ht]
    \centering
    \begin{circuitikz}[scale=0.82, transform shape, line width=0.6pt]
        % --- трубка ---
        \draw[rounded corners=10pt, fill=gray!8] (0,0) rectangle (9,1.8);
        \node[font=\small, anchor=north] at (4.5,-0.05) {};
        % катод
        \draw[line width=2pt] (0.5,0.35) -- (0.5,1.45);
        \node[font=\small, anchor=west] at (0.6,0.4) {К};
        % анод I
        \draw[line width=2pt] (3,1.8) -- (3,1.3);
        \node[font=\small, anchor=north] at (3,1.25) {А-I};
        % анод II
        \draw[line width=2pt] (8.5,0.35) -- (8.5,1.45);
        \node[font=\small, anchor=east] at (8.4,0.4) {А-II};
        % зонды
        \draw[line width=1.4pt] (5.6,0.5) -- (5.6,1.2);
        \draw[line width=1.4pt] (6.0,0.5) -- (6.0,1.2);
        \node[font=\small, anchor=south] at (5.8,1.2) {зонды};

        % --- цепь разряда ---
        \draw (0.5,0.9) -- (-1.5,0.9) -- (-1.5,4.4)
              to[battery1, l_=ВИП] (0.8,4.4)
              to[R, l=$R_{\text{б}}$] (3.2,4.4)
              to[rmeter, t=mA, l=$A_1$] (5.6,4.4) coordinate (top);
        % переключатель П1
        \draw (top) -- (6.6,4.4) node[circ]{} coordinate (sw);
        \draw (sw) -- ++(0.75,0.35) node[ocirc]{};
        \draw (7.5,4.75) -- (10.0,4.75) -- (10.0,0.9) -- (8.5,0.9);
        \draw (7.5,4.05) node[ocirc]{} -- (7.5,3.0) -- (3,3.0) -- (3,1.8);
        \node[font=\small, anchor=south] at (7.0,4.8) {П1};
        % делитель и V1
        \draw (top) -- (5.6,5.6) to[R, l_=$R_1$] (2.2,5.6) coordinate (mid)
              to[R, l_=$R_2$] (-2.8,5.6) -- (-2.8,0.9) -- (-1.5,0.9);
        \draw (mid) -- (2.2,6.8) to[rmeter, t=V, l_=$V_1$] (-2.8,6.8) -- (-2.8,5.6);

        % --- цепь зондов ---
        \draw (5.6,0.5) -- (5.6,-0.9) -- (3.0,-0.9) coordinate (pa);
        \draw (6.0,0.5) -- (6.0,-0.6) -- (8.4,-0.6) -- (8.4,-0.9) coordinate (pb);
        \draw (4.0,-0.9) node[circ]{} -- (4.0,-1.9)
              to[rmeter, t=V, l_=$V_2$] (7.4,-1.9) -- (7.4,-0.9) -- (pb);
        \draw (7.4,-0.9) node[circ]{};
        \draw (pa) to[rmeter, t=\textmu A, l_=$A_2$] (3.0,-3.4) -- (4.3,-3.4);
        \draw (pb) -- (8.4,-3.4) -- (6.7,-3.4);
        \draw[fill=white] (4.3,-3.9) rectangle (6.7,-2.9);
        \node[font=\small] at (5.5,-3.4) {П2};
        \draw (3.0,-5.6) to[pR, n=pot] (8.4,-5.6);
        \node[font=\small] at (5.7,-6.15) {$R$};
        \draw (4.9,-3.9) -- (4.9,-4.6) -| (pot.wiper);
        \draw (6.1,-3.9) -- (6.1,-4.6) -- (8.4,-4.6) -- (8.4,-5.6);
        \draw (3.0,-5.6) -- (3.0,-7.0)
              to[battery1, l_=источник пост. тока] (8.4,-7.0) -- (8.4,-5.6);
    \end{circuitikz}
    \caption{Схема установки для исследования газового разряда
             (упрощённо; по рис.~1 описания работы)}
    \label{fig:scheme}
\end{figure}

\paragraph{ВАХ разряда.} Переключатель П1 устанавливается в положение
«Анод-I». Плавно увеличивая напряжение ВИП, фиксируют напряжение
зажигания $U_{\text{заж}}$ (показание $V_1$ непосредственно перед
зажиганием). Затем измеряют $I_{\text{р}}(U_{\text{р}})$ при увеличении и
при уменьшении тока в допустимом диапазоне $I_{\text{р}} \leqslant \qq$~мА
и отмечают напряжение гашения разряда.

\paragraph{Зондовые характеристики.} Напряжение ВИП сбрасывается до нуля,
П1 переводится в положение «Анод-II»: разряд горит между катодом и
анодом-II, и двойной зонд оказывается в плазме положительного столба.
Устанавливается ток разряда $I_{\text{р}} \approx 5$~мА, потенциометром $R$
--- максимальное напряжение на зондах $U_{\text{з}}^{\max} \approx 25$~В.
ВАХ зонда $I_{\text{з}}(U_{\text{з}})$ снимается в диапазоне
$-U_{\text{з}}^{\max}\ldots U_{\text{з}}^{\max}$; при $I_{\text{з}} = 0$
полярность меняется переключателем П2, вблизи $U_{\text{з}} = 0$ точки
берутся чаще. Ток разряда в процессе измерения поддерживается постоянным.
Измерения повторяются ещё при 3--4 значениях $I_{\text{р}}$.

\paragraph{Обработка.} Зондовые характеристики центрируются (сдвиг по току
--- на полусумму пересечений асимптот с осью $U = 0$, по напряжению --- так,
чтобы $I(0) = 0$). Ток насыщения $I_{i\text{н}}$ определяется как
полуразность пересечений асимптот (МНК по точкам с
$|U| > U_{\text{ас}}$) с осью $U = 0$, наклон $dI/dU|_{0}$ --- МНК по точкам с
$|U| < U_{\text{лин}}$. Для контроля характеристика аппроксимируется
функцией $I = I_{i\text{н}}\operatorname{th}\bigl(eU/2k_{\text{Б}}T_e\bigr) + AU$.
Погрешности $I_{i\text{н}}$ и $dI/dU$ --- случайные погрешности МНК,
погрешности косвенных величин находятся по правилам переноса.

\clearpage
\section{Результаты измерений и обработка данных}

\subsection{Параметры установки}
\begin{itemize}
    \item диаметр и длина зонда: $d = \qq$~мм, $l = \qq$~мм,
          площадь $S = \pi d l = \qq$~мм$^2$;
    \item коэффициент делителя $(R_1+R_2)/R_2 = \qq$;
    \item балластное сопротивление $R_{\text{б}} = \qq$~кОм;
    \item давление неона $P \approx 2$~торр, $T \approx 300$~К,
          $n_0 = P/k_{\text{Б}}T \approx 6{,}4\cdot10^{16}$~см$^{-3}$;
    \item границы участков обработки: $U_{\text{ас}} = \qq$~В,
          $U_{\text{лин}} = \qq$~В.
\end{itemize}
Все величины дублируются в \filename{data/raw/setup_params.csv} и
используются скриптами обработки.

% ---------------------------------------------------------------------------
\subsection{Вольт-амперная характеристика разряда}
Напряжение зажигания разряда $U_{\text{заж}} = \qq$~В, напряжение
гашения $U_{\text{гаш}} = \qq$~В. Результаты измерения ВАХ приведены в
\cref{tab:discharge} и на \cref{fig:discharge}.

\begin{table}[!ht]
    \centering
    \caption{ВАХ тлеющего разряда (П1 в положении «Анод-I»);
             $R_{\text{диф}}$ --- локальная производная $dU/dI$}
    \label{tab:discharge}
    \IfFileExists{\proc{discharge_iv_processed.csv}}{%
    \csvreader[
        tabularray = {
            colspec = {|c|c|c|c|c|},
            hlines, vlines,
            row{1} = {font=\bfseries},
        },
        table head = {
            ветвь & $U_{V_1}$, В & $U_{\text{р}}$, В & $I_{\text{р}}$, мА
            & $R_{\text{диф}}$, кОм \\
        }
    ]%
    {\proc{discharge_iv_processed.csv}}%
    {branch=\cBr, U_V1_V=\cUv, U_r_V=\cUr, I_mA=\cI, R_diff_kOhm=\cR}%
    {%
        \ifcsvstrcmp{\cBr}{up}{$\uparrow$}{$\downarrow$}
        & $\pnum{\cUv}$ & $\pnum{\cUr}$ & $\pnum{\cI}$ & $\pnum{\cR}$
    }%
    }{%
    \small
    \BlankTable{Q[c,25mm]|Q[c,25mm]||Q[c,25mm]|Q[c,25mm]|Q[c,25mm]}%
        {$U_{V_1}$, В $\uparrow$ & $I_{\text{р}}$, мА $\uparrow$ &
         $U_{V_1}$, В $\downarrow$ & $I_{\text{р}}$, мА $\downarrow$ &
         примечание}{5}{16}%
    }
\end{table}

\begin{figure}[!ht]
    \centering
    \IfFileExists{\proc{discharge_iv_up.csv}}{%
    \begin{tikzpicture}
        \begin{axis}[
            width=0.85\textwidth, height=0.55\textwidth,
            xlabel={$U_{\text{р}}$, В},
            ylabel={$I_{\text{р}}$, мА},
            ymin=0,
            grid=both, minor tick num=1,
            legend pos=north east,
            /pgf/number format/use comma,
        ]
            \addplot[only marks, mark=*, mark size=1.8pt]
                table[col sep=comma, x=U_r_V, y=I_mA] {\proc{discharge_iv_up.csv}};
            \addlegendentry{увеличение тока}
            \IfFileExists{\proc{discharge_iv_down.csv}}{%
            \addplot[only marks, mark=o, mark size=2.2pt]
                table[col sep=comma, x=U_r_V, y=I_mA] {\proc{discharge_iv_down.csv}};
            \addlegendentry{уменьшение тока}
            }{}
        \end{axis}
    \end{tikzpicture}%
    }{\labimage{discharge_iv.pdf}{0.8}}
    \caption{Вольт-амперная характеристика тлеющего разряда в неоне}
    \label{fig:discharge}
\end{figure}

Наклон ВАХ в координатах $U_{\text{р}}(I_{\text{р}})$ (МНК):
\begin{equation}
    R_{\text{диф}} = \frac{dU_{\text{р}}}{dI_{\text{р}}} = \qq \pm \qq~\text{кОм},
\end{equation}
максимальное по модулю дифференциальное сопротивление
$R_{\text{диф}}^{\max} = \qq$~кОм (при $I_{\text{р}} = \qq$~мА).

Сравнение с рис.~6 Приложения к разделу V: полученный участок ВАХ
соответствует \fillme{} (поднормальному / нормальному / аномальному
тлеющему разряду), поскольку \fillme.

% ---------------------------------------------------------------------------
\clearpage
\subsection{Зондовые характеристики}
Зондовые характеристики сняты при токах разряда
$I_{\text{р}} = \qq,\ \qq,\ \qq,\ \qq$~мА. Семейство центрированных
характеристик показано на \cref{fig:probe_family}, каждая характеристика
с построенными асимптотами и касательной в нуле --- на
\cref{fig:probe_single}.

\begin{figure}[!ht]
    \centering
    \IfFileExists{\proc{probe_family.tex}}{%
    \begin{tikzpicture}
        \begin{axis}[
            width=0.9\textwidth, height=0.62\textwidth,
            xlabel={$U$, В},
            ylabel={$I$, мкА},
            axis lines=middle,
            grid=both, minor tick num=1,
            legend pos=south east,
            legend style={font=\small},
            xlabel style={at={(ticklabel* cs:1)}, anchor=north west},
            ylabel style={at={(ticklabel* cs:1)}, anchor=south west},
            cycle list name=exotic,
        ]
            \input{\proc{probe_family.tex}}
        \end{axis}
    \end{tikzpicture}%
    }{\labimage{probe_family.pdf}{0.85}}
    \caption{Семейство центрированных ВАХ двойного зонда при различных
             токах разряда}
    \label{fig:probe_family}
\end{figure}

\begin{figure}[!ht]
    \centering
    \begin{subfigure}{0.48\textwidth}
        \centering
        \labimage{probe_1.pdf}{1.0}
        \caption{$I_{\text{р}} = \qq$~мА}
    \end{subfigure}
    \hfill
    \begin{subfigure}{0.48\textwidth}
        \centering
        \labimage{probe_2.pdf}{1.0}
        \caption{$I_{\text{р}} = \qq$~мА}
    \end{subfigure}

    \vspace{4mm}
    \begin{subfigure}{0.48\textwidth}
        \centering
        \labimage{probe_3.pdf}{1.0}
        \caption{$I_{\text{р}} = \qq$~мА}
    \end{subfigure}
    \hfill
    \begin{subfigure}{0.48\textwidth}
        \centering
        \labimage{probe_4.pdf}{1.0}
        \caption{$I_{\text{р}} = \qq$~мА}
    \end{subfigure}
    \caption{ВАХ двойного зонда: асимптоты (пунктир) и касательная в
             начале координат; квадратами отмечены $\pm I_{i\text{н}}$
             (графики \filename{scripts/plot_data.py})}
    \label{fig:probe_single}
\end{figure}

По каждой характеристике определены ионный ток насыщения $I_{i\text{н}}$ и
наклон $dI/dU|_{U=0}$, по формуле \eqref{eq:Te} рассчитана электронная
температура (\cref{tab:probe}). В последнем столбце --- $k_{\text{Б}}T_e$,
полученная аппроксимацией всей характеристики функцией
\eqref{eq:tanh} с линейной добавкой.

\begin{table}[!ht]
    \centering
    \caption{Результаты обработки зондовых характеристик}
    \label{tab:probe}
    \IfFileExists{\proc{probe_results.csv}}{%
    \csvreader[
        tabularray = {
            colspec = {|c|c|c|c|c|},
            hlines, vlines,
            row{1} = {font=\bfseries},
        },
        table head = {
            $I_{\text{р}}$, мА & $I_{i\text{н}}$, мкА
            & $dI/dU|_0$, мкА/В & $k_{\text{Б}}T_e$, эВ
            & $k_{\text{Б}}T_e^{\,\text{th}}$, эВ \\
        }
    ]%
    {\proc{probe_results.csv}}%
    {Ip_mA=\cIp, Iin_uA=\cIin, sIin_uA=\csIin, dIdU_uA_V=\cK, sdIdU_uA_V=\csK,
     Te_eV=\cTe, sTe_eV=\csTe, Te_fit_eV=\cTef}%
    {%
        $\pnum{\cIp}$ & $\pnum{\cIin} \pm \pnum{\csIin}$
        & $\pnum{\cK} \pm \pnum{\csK}$ & $\pnum{\cTe} \pm \pnum{\csTe}$
        & $\pnum{\cTef}$
    }%
    }{%
    \BlankTable{Q[c,20mm]|Q[c,28mm]|Q[c,28mm]|Q[c,28mm]|Q[c,28mm]}%
        {$I_{\text{р}}$, мА & $I_{i\text{н}}$, мкА & $dI/dU|_0$, мкА/В
         & $k_{\text{Б}}T_e$, эВ & $k_{\text{Б}}T_e^{\,\text{th}}$, эВ}{5}{4}%
    }
\end{table}

% ---------------------------------------------------------------------------
\subsection{Параметры плазмы}
По формулам \eqref{eq:si1}--\eqref{eq:si3} рассчитаны концентрация
электронов, плазменная частота, электронная поляризационная длина,
радиус экранирования (с учётом $T_i \approx 300$~К $\ll T_e$), число
частиц в дебаевской сфере и степень ионизации (\cref{tab:plasma}).

\begin{table}[!ht]
    \centering
    \caption{Параметры плазмы положительного столба при различных токах
             разряда}
    \label{tab:plasma}
    \small
    \IfFileExists{\proc{plasma_params.csv}}{%
    \csvreader[
        tabularray = {
            colspec = {|c|c|c|c|c|c|c|c|c|},
            hlines, vlines,
            row{1} = {font=\bfseries},
        },
        table head = {
            $I_{\text{р}}$, мА & $k_{\text{Б}}T_e$, эВ & $T_e$, К
            & $n_e$, $10^{10}$~см$^{-3}$ & $\omega_p$, $10^{9}$~с$^{-1}$
            & $r_{De}$, мкм & $r_D$, мкм & $N_D$ & $\alpha$, $10^{-7}$ \\
        }
    ]%
    {\proc{plasma_params.csv}}%
    {Ip_mA=\cIp, Te_eV=\cTe, sTe_eV=\csTe, Te_K=\cTK, n_1e10_cm3=\cn,
     sn_1e10_cm3=\csn, wp_1e9_rad_s=\cwp, rDe_um=\crDe, rD_um=\crD, ND=\cND,
     alpha_1e7=\cal}%
    {%
        $\pnum{\cIp}$ & $\pnum{\cTe}$ & $\pnum{\cTK}$
        & $\pnum{\cn} \pm \pnum{\csn}$ & $\pnum{\cwp}$ & $\pnum{\crDe}$
        & $\pnum{\crD}$ & $\pnum{\cND}$ & $\pnum{\cal}$
    }%
    }{%
    \BlankTable{Q[c,13mm]|Q[c,14mm]|Q[c,16mm]|Q[c,17mm]|Q[c,16mm]|Q[c,13mm]|Q[c,13mm]|Q[c,11mm]|Q[c,13mm]}%
        {$I_{\text{р}}$, мА & $k_{\text{Б}}T_e$, эВ & $T_e$, К
         & $n_e$, $10^{10}$ см$^{-3}$ & $\omega_p$, $10^{9}$ с$^{-1}$
         & $r_{De}$, мкм & $r_D$, мкм & $N_D$ & $\alpha$, $10^{-7}$}{9}{4}%
    }
\end{table}

\begin{itemize}
    \item \textbf{Квазинейтральность.} $r_D \approx \qq$~мкм и
          $r_{De} \approx \qq$~мкм много меньше радиуса трубки
          $a \approx \qq$~мм, следовательно \fillme.
    \item \textbf{Идеальность.} $N_D \approx \qq \gg 1$ (или нет), поэтому
          плазма \fillme.
    \item \textbf{Равновесность.} $T_e \approx \qq$~К $\gg T_i \approx 300$~К
          --- плазма \fillme.
    \item \textbf{Степень ионизации} $\alpha \sim \qq$ --- плазма
          слабоионизованная (\fillme).
\end{itemize}

\begin{figure}[!ht]
    \centering
    \IfFileExists{\proc{plasma_params.csv}}{%
    \begin{subfigure}{0.49\textwidth}
        \centering
        \begin{tikzpicture}
            \begin{axis}[
                width=\linewidth, height=0.85\linewidth,
                xlabel={$I_{\text{р}}$, мА}, ylabel={$k_{\text{Б}}T_e$, эВ},
                xmin=0, ymin=0, grid=both,
            ]
                \addplot[only marks, mark=*, mark size=1.8pt,
                         error bars/.cd, y dir=both, y explicit]
                    table[col sep=comma, x=Ip_mA, y=Te_eV, y error=sTe_eV]
                    {\proc{plasma_params.csv}};
            \end{axis}
        \end{tikzpicture}
        \caption{$T_e(I_{\text{р}})$}
    \end{subfigure}
    \hfill
    \begin{subfigure}{0.49\textwidth}
        \centering
        \begin{tikzpicture}
            \begin{axis}[
                width=\linewidth, height=0.85\linewidth,
                xlabel={$I_{\text{р}}$, мА}, ylabel={$n_e$, $10^{10}$~см$^{-3}$},
                xmin=0, ymin=0, grid=both,
            ]
                \addplot[only marks, mark=*, mark size=1.8pt,
                         error bars/.cd, y dir=both, y explicit]
                    table[col sep=comma, x=Ip_mA, y=n_1e10_cm3, y error=sn_1e10_cm3]
                    {\proc{plasma_params.csv}};
            \end{axis}
        \end{tikzpicture}
        \caption{$n_e(I_{\text{р}})$}
    \end{subfigure}%
    }{%
    \begin{subfigure}{0.48\textwidth}
        \centering\labimage{Te_vs_Ip.pdf}{1.0}\caption{$T_e(I_{\text{р}})$}
    \end{subfigure}
    \hfill
    \begin{subfigure}{0.48\textwidth}
        \centering\labimage{ne_vs_Ip.pdf}{1.0}\caption{$n_e(I_{\text{р}})$}
    \end{subfigure}%
    }
    \caption{Зависимость электронной температуры и концентрации электронов
             от тока разряда}
    \label{fig:Te_ne}
\end{figure}

Анализ зависимостей: \fillme{} (ожидается рост $n_e$ примерно
пропорционально $I_{\text{р}}$, поскольку плотность тока в
положительном столбе $j \approx e n_e \mu_e E$; температура электронов
определяется балансом ионизации и диффузионных потерь на стенки и
слабо зависит от тока).

\subsection{Оценка погрешностей}
\begin{itemize}
    \item Приборные погрешности: $A_1$ --- \qq, $V_1$ --- \qq,
          $A_2$ --- \qq, $V_2$ --- \qq.
    \item Погрешности $I_{i\text{н}}$ и $dI/dU|_0$ определяются по
          разбросу точек (МНК); относительная погрешность $T_e$:
          $\varepsilon_{T} = \sqrt{\varepsilon_{I}^2 + \varepsilon_{k}^2}
          \approx \qq$~\%.
    \item Относительная погрешность концентрации
          $\varepsilon_n = \sqrt{\varepsilon_I^2 + \varepsilon_T^2/4
          + \varepsilon_S^2} \approx \qq$~\%. Систематическую
          погрешность вносят полуэмпирический коэффициент $0{,}4$ в формуле
          Бома и неопределённость площади зонда из-за дебаевского слоя
          (\fillme).
    \item Результат $T_e$ зависит от выбора участков для асимптот и
          касательной: при изменении $U_{\text{лин}}$ в пределах \qq{}
          $T_e$ меняется на \qq~\%. Сравнение с аппроксимацией
          функцией th: \fillme.
\end{itemize}

\section{Обсуждение результатов}
\begin{itemize}
    \item ВАХ разряда: \fillme.
    \item Зондовые характеристики имеют вид, предсказанный
          \eqref{eq:tanh}; наклон асимптот \fillme.
    \item Полученные $k_{\text{Б}}T_e \approx \qq$~эВ и
          $n_e \sim \qq$~см$^{-3}$ сравнить с типичными для тлеющего
          разряда значениями ($T_e \sim 1$--$5$~эВ,
          $n_e \sim 10^{9}$--$10^{11}$~см$^{-3}$): \fillme.
    \item Зависимости $T_e(I_{\text{р}})$ и $n_e(I_{\text{р}})$: \fillme.
\end{itemize}

\section{Заключение}
\begin{itemize}
    \item Измерена ВАХ тлеющего разряда в неоне: $U_{\text{заж}} = \qq$~В,
          максимальное дифференциальное сопротивление
          $R_{\text{диф}} = \qq$~кОм; исследованный участок соответствует
          \fillme.
    \item Методом двойного зонда при $I_{\text{р}} = \qq\div\qq$~мА
          получены $k_{\text{Б}}T_e = \qq\div\qq$~эВ
          ($T_e \approx \qq$~К) и $n_e = \qq\div\qq$~см$^{-3}$.
    \item Плазменная частота $\omega_p \approx \qq$~с$^{-1}$, поляризационная
          длина $r_{De} \approx \qq$~мкм, радиус экранирования
          $r_D \approx \qq$~мкм, $N_D \approx \qq$, степень ионизации
          $\alpha \approx \qq$.
    \item Плазма положительного столба квазинейтральна (\fillme),
          идеальна (\fillme), неравновесна ($T_e \gg T_i$) и слабо
          ионизована.
\end{itemize}

\end{document}
